% ==============================================================================
% INTRODUCTION GÉNÉRALE — NG-IDPS
% Système Autonome de Détection et Prévention d'Intrusions Nouvelle Génération
% Master Professionnel : Computer Engineering — Expert Réseaux, ISIMa
% ==============================================================================

% \chapter*{Introduction Générale}
\label{chap:intro}
% ==============================================================================
% INTRODUCTION GÉNÉRALE
% ==============================================================================
À l'ère de la transformation numérique, l'interconnexion croissante des systèmes d'information et la dépendance des organisations vis-à-vis de leurs infrastructures réseau marquent une rupture significative avec les pratiques traditionnelles de la sécurité informatique. Le paysage des menaces évolue en parallèle : les cyberattaques deviennent plus variées, plus automatisées et plus difficiles à caractériser à l'avance. La défense moderne ne peut plus reposer uniquement sur des règles prédéfinies ; elle s'appuie désormais sur la complémentarité de plusieurs mécanismes : détection par signatures, analyse comportementale fondée sur l'apprentissage automatique, corrélation de sources multiples, renseignement sur les menaces et automatisation de la réponse. Cette convergence redéfinit les standards de réactivité, de résilience et d'efficacité des centres opérationnels de sécurité (\textit{Security Operations Center}, SOC).

Dans un pays en développement comme la Tunisie, l'adoption de ces approches reste encore limitée. De nombreuses organisations disposent de ressources humaines et financières restreintes et ne peuvent pas assurer une supervision permanente de leurs réseaux. Elles s'appuient souvent sur des mécanismes de protection statiques, les exposant à des délais de détection importants, à des alertes difficiles à exploiter et à des attaques qu'aucune règle connue ne décrit encore. Par ailleurs, les plateformes commerciales intégrées reposent fréquemment sur des licences ou des services dans le nuage qui ne conviennent pas toujours aux organisations soucieuses de conserver la maîtrise de leurs données. Cette situation freine la mise en place d'une stratégie de cyberdéfense proactive et autonome.

C'est dans cette optique que s'inscrit le présent projet de fin d'études, réalisé à l'Institut Supérieur d'Informatique de Mahdia (ISIMa), dans le cadre du Mastère Professionnel en Computer Engineering, spécialité Expert Réseaux. L'objectif principal du projet est de concevoir, de réaliser et de valider un système de détection et de prévention des intrusions réseau de nouvelle génération, désigné \textbf{NG-IDPS} (\textit{Next-Generation Intrusion Detection and Prevention System}), capable de fonctionner de manière autonome dans un environnement local et isolé. Ce système vise à combiner plusieurs mécanismes de détection complémentaires, à analyser et à corréler les événements de sécurité, et à déclencher des actions de réponse maîtrisées, tout en restant accessible à des organisations disposant de moyens limités.

Le présent mémoire retrace les différentes étapes de conception, de développement et de validation de ce projet. Il est structuré en quatre chapitres, suivis d'une conclusion générale :

\begin{itemize}
    \item \textbf{Le premier chapitre} pose le cadre du projet en présentant son contexte, sa problématique et ses objectifs, en étudiant les solutions existantes et en exposant la méthodologie de travail retenue.

    \item \textbf{Le deuxième chapitre} est consacré à la phase préparatoire : l'analyse et la spécification des besoins fonctionnels et non fonctionnels, la planification selon la méthodologie agile Scrum, la justification des choix technologiques et la présentation de l'architecture générale.

    \item \textbf{Le troisième chapitre} traite de la première \textit{Release}, dédiée à la détection : l'acquisition du trafic réseau, l'acheminement des données et la détection comportementale, ainsi que leur validation expérimentale.

    \item \textbf{Le quatrième chapitre} détaille la seconde \textit{Release}, dédiée à la prévention et à la réponse : l'anticipation des menaces par le renseignement, la sensibilisation des utilisateurs, l'orchestration de la réponse automatisée et la supervision par un tableau de bord dédié.

    \item \textbf{La conclusion générale} dresse le bilan du travail réalisé, en discute les limites et propose des perspectives d'évolution.
\end{itemize}

À travers ce travail, nous avons cherché à apporter une réponse concrète aux besoins de protection des infrastructures réseau des organisations tunisiennes, en proposant une solution rigoureuse, maîtrisable localement et intégrable à des environnements existants.